\chapter{Tema d'esame di un appello precedente}
\label{cap:tema-precedente}

\begin{nota}[Fonte]
Messaggio nel gruppo Telegram del corso (``per i posteri''), riportato testualmente. Il testo delle
opzioni delle domande chiuse non è disponibile: per ciascuna si indica cosa va saputo. Le domande
chiuse ricostruite si trovano nel quiz (Capitolo~\ref{cap:quiz}).
\end{nota}

\begin{enumerate}
  \item \emph{GDPR legal basis (select the WRONG answer)}
  \item \emph{NIS2 incident definition}
  \item \emph{ENISA: what is it, what role does it have}
  \item \emph{DORA strategies (select the WRONG answer)}
  \item \emph{Open question: a hospital with some IT service providers for storage (cloud) and security.
    There is a data breach/incident.}
    \begin{enumerate}[label=\Alph*)]
      \item \emph{What regulation is involved?}
      \item \emph{According to the regulation, what does the hospital have to do?}
      \item \emph{Regardless of the incident, what does the hospital have to do regarding the
        management of the supply chain?}
    \end{enumerate}
\end{enumerate}

% ---------------------------------------------------------------------
\section{Domande chiuse: cosa sapere}

\paragraph{1. Basi giuridiche GDPR (art.~6).} Sono sei, tutte di pari valore: consenso, esecuzione di
un contratto, obbligo legale, interessi vitali, interesse pubblico o pubblici poteri, legittimo
interesse. Affermazioni errate tipiche:
\begin{itemize}
  \item ``il consenso è sempre necessario'' o ``prevale sulle altre basi'';
  \item ``un'autorità pubblica può fondarsi sul legittimo interesse per i propri compiti'' (escluso
    dall'art.~6(1));
  \item ``la base contrattuale copre anche il marketing'';
  \item ``si può condividere i dati con terzi per il loro marketing in base al legittimo interesse''
    (laboratorio GDPR: serve il consenso).
\end{itemize}

\paragraph{2. Definizione di incidente NIS2 (art.~6).} Evento che compromette la
\textbf{disponibilità, l'autenticità, l'integrità o la riservatezza} di dati conservati, trasmessi o
elaborati, o dei servizi offerti dai sistemi informativi e di rete o accessibili attraverso di essi.
Va distinto dall'\textbf{incidente significativo} (art.~23(3)): quello che ha causato o può causare
una grave perturbazione operativa dei servizi o perdite finanziarie, oppure danni materiali o
immateriali considerevoli ad altre persone. Il ``può causare'' è giuridicamente decisivo.
Attenzione: NIS2 include l'\emph{autenticità}, la definizione GDPR di data breach no.

\paragraph{3. ENISA.} Agenzia dell'UE per la cybersicurezza. Con il Cybersecurity Act ha ottenuto un
\textbf{mandato permanente} e un ruolo attivo: supporto alle politiche UE e agli Stati membri,
cooperazione operativa (rete dei CSIRT), situational awareness, preparazione degli schemi europei di
certificazione su richiesta della Commissione, sensibilizzazione e formazione, gestione della
\textbf{Single Reporting Platform} del CRA. \textbf{Non} irroga sanzioni, \textbf{non} è un CSIRT
nazionale e \textbf{non} rilascia certificati alle imprese.

\paragraph{4. Strategie DORA.} La strategia di resilienza operativa digitale comprende: tolleranza al
rischio ICT, obiettivi, architettura, rilevazione, test; \textbf{strategia multi-vendor}; gestione e
conformità dell'outsourcing ICT; \textbf{strategia di comunicazione} per gli incidenti ICT. È definita
e approvata dall'\textbf{organo di gestione} dell'entità finanziaria. Affermazioni errate tipiche: che
la approvi il fornitore ICT, che sia identica al risk assessment GDPR, che DORA imponga una checklist
fissa (vale invece la proporzionalità).

% ---------------------------------------------------------------------
\section{Domanda aperta: risposta modello}

\begin{soluzione}[A) Quali norme sono coinvolte]
\begin{itemize}
  \item \textbf{GDPR}. L'ospedale tratta dati personali, anche \textbf{sanitari} (categoria
    particolare, art.~9), ed è \textbf{titolare}. Il fornitore cloud e il fornitore di sicurezza sono
    \textbf{responsabili del trattamento} (art.~28), ed eventuali loro fornitori sono sub-responsabili.
    Se sono coinvolti dati personali, l'incidente è un \textbf{data breach} (art.~4(12)).
  \item \textbf{NIS2}. La sanità è un settore dell'Allegato~I: l'ospedale è \textbf{soggetto
    essenziale} se grande impresa, \textbf{importante} se media (art.~3). In Italia la norma è il
    D.Lgs.~138/2024, con l'ACN come autorità e il CSIRT Italia per le notifiche.
  \item \textbf{I fornitori} sono a loro volta soggetti NIS2: i servizi di cloud computing rientrano
    nelle infrastrutture digitali e i servizi di sicurezza gestiti nella gestione dei servizi ICT
    (B2B), entrambi nell'Allegato~I. Hanno obblighi di notifica propri, distinti da quelli
    dell'ospedale.
  \item \textbf{Perché non le altre}: DORA no (l'ospedale non è un'entità finanziaria); CRA solo se
    l'incidente nasce da una vulnerabilità attivamente sfruttata in un prodotto con elementi digitali,
    e in quel caso l'obbligo è del fabbricante, non dell'ospedale; i dispositivi medici sono soggetti
    all'MDR ed esclusi dal CRA.
\end{itemize}
\end{soluzione}

\begin{soluzione}[B) Cosa deve fare l'ospedale secondo le norme]
\textbf{Gestione immediata}
\begin{itemize}
  \item Attivare la procedura (idealmente integrata GDPR--NIS2), contenere l'incidente insieme ai
    fornitori, preservare le evidenze (log, immagini forensi, catena di custodia) e informare l'organo
    di gestione.
  \item Ottenere dal fornitore le informazioni necessarie: il responsabile deve informare il titolare
    \textbf{senza ingiustificato ritardo} (art.~33(2)).
  \item Fissare e documentare il momento di consapevolezza per ciascuna norma.
\end{itemize}
\textbf{GDPR}
\begin{itemize}
  \item Valutare la severità per gli interessati (ENISA). Con dati sanitari il DPC è 4, quindi la
    severità è quasi certamente alta o molto alta.
  \item \textbf{Art.~33}: notificare il Garante entro \textbf{72 ore} dalla consapevolezza, anche in
    forma preliminare, completando in seguito.
  \item \textbf{Art.~34}: comunicare ai pazienti, in linguaggio chiaro, se il rischio è elevato (con
    dati sanitari tipicamente sì). La comunicazione deve contenere natura della violazione, contatti
    del DPO, conseguenze probabili, misure adottate e consigli agli interessati.
  \item Annotare la violazione nel \textbf{registro dei data breach} (art.~33(5)), anche se non si notifica.
\end{itemize}
\textbf{NIS2} (se l'incidente è significativo)
\begin{itemize}
  \item \textbf{Early warning} al CSIRT entro \textbf{24 ore}, indicando se si sospetta un atto
    illecito o un impatto transfrontaliero.
  \item \textbf{Notifica dell'incidente} entro \textbf{72 ore}, con valutazione iniziale di gravità,
    impatto e indicatori di compromissione; \textbf{relazione intermedia} su richiesta.
  \item \textbf{Relazione finale} entro \textbf{un mese} dalla notifica: descrizione, causa, misure,
    impatti transfrontalieri.
  \item Informare, se opportuno, i destinatari dei servizi impattati.
\end{itemize}
\textbf{Inoltre}
\begin{itemize}
  \item Le notifiche al Garante e al CSIRT sono \textbf{separate e parallele}: l'una non sostituisce l'altra.
  \item Valutare la denuncia alla Polizia Postale (non obbligatoria, ma la sua assenza può essere
    valutata negativamente).
  \item Allineare registro degli incidenti e registro dei data breach; far approvare le misure di
    remediation dall'organo di gestione; aggiornare risk assessment e DPIA.
\end{itemize}
\end{soluzione}

\begin{soluzione}[C) Supply chain: obblighi a prescindere dall'incidente]
\begin{itemize}
  \item \textbf{NIS2, art.~21(2)(d)}: la sicurezza della catena di approvvigionamento fa parte delle
    misure obbligatorie di gestione del rischio. Secondo l'\textbf{art.~21(3)} vanno considerate le
    vulnerabilità specifiche di ciascun fornitore, la qualità dei loro prodotti e delle loro pratiche di
    sicurezza e le procedure di sviluppo sicuro.
  \item In pratica (linee guida ACN): inventario dei fornitori e dei servizi critici; mappatura delle
    dipendenze e dei rischi a cascata; requisiti minimi di sicurezza nei contratti; SLA e tempi di
    ripristino; rivalutazione periodica; supervisione e approvazione da parte dell'organo di gestione
    (art.~20, responsabilità non delegabile).
  \item \textbf{GDPR, art.~28}: scegliere solo responsabili che offrano \textbf{garanzie sufficienti};
    stipulare un contratto (DPA) con il contenuto dell'art.~28(3), cioè istruzioni documentate,
    riservatezza, misure dell'art.~32, regole sui sub-responsabili, assistenza al titolare (anche per
    gli artt.~33--36), cancellazione o restituzione dei dati, diritto di audit.
  \item Prevedere una \textbf{clausola di notifica dei breach} con tempi stretti (es.\ 24 ore), così da
    poter rispettare le 72 ore.
  \item \textbf{Localizzazione dei dati}: se il cloud tratta dati fuori dall'UE servono gli strumenti
    del Capo~V (decisione di adeguatezza oppure SCC con TIA).
  \item Clausole su cooperazione negli incidenti, continuità operativa, backup, \textbf{exit strategy}
    e reversibilità.
  \item Le certificazioni dei fornitori (ISO 27001, 27017, 27018) sono utili come evidenza ma
    \textbf{non bastano} da sole: le verifiche vanno documentate (accountability).
\end{itemize}
\end{soluzione}
